\section{Problem Formulation}
\label{sec:problem-formulation}
\paragraph{Sequential future-state forecasting.}
Given an observed state context $x_{a:m}$, where
$x_k\in\mathcal X\subseteq\R^d$ and $a\leq m<T$, we seek to model its future
continuation $x_{m+1:T}$.  We learn a conditional generative model
$Q_\theta(\cdot\mid x_{a:m})$ and draw forecasts
$\widehat x_{m+1:T}\sim Q_\theta(\cdot\mid x_{a:m})$.  This formulation
covers three settings, distinguished by the predictive target and the role of
forecast accuracy at inference:

\par\smallskip
\noindent\textbf{Fully observed (deterministic).}
When the system state is fully observed and the dynamics are deterministic,
the context uniquely determines the future,
$x_{m+1:T}=F_\star(x_{a:m})$.  For the point-valued inference goal considered
here, the predictive mean should recover this future,
$\E_{Q_\theta(\cdot\mid x_{a:m})}[\widehat x_{m+1:T}]=x_{m+1:T}$.

\par\smallskip
\noindent\textbf{Partially observed (stochastic).}
When relevant variables are unobserved or the dynamics are intrinsically
stochastic, the same observed context can admit multiple futures, described by
$p_\star(\cdot\mid x_{a:m})$.  Forecast accuracy is therefore distributional:
$Q_\theta$ should represent the full conditional law, including its
variability and multimodality, rather than only its conditional mean.

\par\smallskip
\noindent\textbf{World modeling.}
In world modeling, $a_k$ drives the transition $x_k\to x_{k+1}$, and the
controlled predictor conditions on the observed actions $a_{a:m-1}$ and a
candidate future action sequence $a_{m:T-1}$, jointly denoted by
$a_{a:T-1}:=(a_{a:m-1},a_{m:T-1})$.  Given a goal $x_g$ and terminal cost
$d_g$, a policy $\pi$ induces the model-imagined terminal law
$Q_{\theta,T}^{\pi}$ and is evaluated by its expected terminal cost.  If
$Q_\theta$ matches the environment future-sequence law for every admissible
future action sequence, then the terminal laws, and hence the expected cost of
every policy, also agree.  The model-based and environment objectives
therefore have the same minimizers
(Eq.~\ref{eq:problem_exact_planning_consistency}).  For squared terminal
distance, this conclusion concerns expected goal cost.
Table~\ref{tab:problem_settings} summarizes the three settings.

\begin{table*}[t]
  \caption{Ground truth and inference goal for the three sequential
  prediction settings.  World-model inference
  first selects a goal-conditioned planning policy by minimizing expected
  model-imagined terminal goal cost, and then samples action trajectories from
  that policy; real success is evaluated using the environment rollout law.}
  \label{tab:problem_settings}
  \centering
  \small
  \setlength{\tabcolsep}{6pt}
  \renewcommand{\arraystretch}{1.2}
  \resizebox{\linewidth}{!}{%
  \begin{tabular}{@{}lll@{}}
    \toprule
    Setting
      & Ground truth
      & Inference goal \\
    \midrule
    Deterministic
      & $x_{m+1:T}=F_\star(x_{a:m})$
      & $\displaystyle
         \E_{Q_\theta(\cdot\mid x_{a:m})}
         [\widehat x_{m+1:T}]=x_{m+1:T}$ \\
    Stochastic
      & $x_{m+1:T}\sim p_\star(\cdot\mid x_{a:m})$
      & $\widehat x_{m+1:T}
         \sim Q_\theta(\cdot\mid x_{a:m})$ \\
    World model
      & $x^{\mathrm{env}}_{m+1:T}\sim
         p_{\mathrm{env}}(\cdot\mid x_{a:m},a_{a:T-1})$
      & $\displaystyle
         \pi_\theta^\star\in
         \argmin_{\pi}
         \E_{\widehat x_T\sim Q_{\theta,T}^{\pi}}
         \!\left[d_g(\widehat x_T,x_g)\right],\;
         a_{m:T-1}\sim\pi_\theta^\star$ \\
    \bottomrule
  \end{tabular}}
\end{table*}
